There are two main parts of the work in this thesis.
Each of them focus on one specific \emph{quantitative-reachability} property,
namely the \emph{adaptivity} of an adadptive data analysis program in \redd{PART \romannum{1}}
and \emph{the number of times that each program location is visited during the program execution} in \redd{PART \romannum{2}}.
% The first part focuses on formalizing and analyzing the \emph{adaptivity} -- a quantitative reachability property -- for the programs in the data analysis area,
% through semantics-based formalization and static program analysis techniques.
% The second part focus on computing the precise reachability-bound for every program control point.
\begin{enumerate}

 \item \redd{PART \romannum{1} \quad PROGRAM ANALYSIS FRAMEWORK FOR ADAPTIVE DATA ANALYSIS}
 shows the work on studying and formalizing the \emph{adaptivity} of the adaptive data analysis program. 

 \paragraph{Chapter~\ref{sec:adapt-intro}} has four sections.
 
 In Section~\ref{sec:adapt-background}, we first introduce the background of
 the adaptive data analysis in theory and algorithm research areas.
 %
 Then in Section~\ref{sec:adapt-motivation}, we discuss in detail some results
 about \emph{adaptivity} of the adaptive data analysis from the theory areas and the challenges and motivations in analyzing
 this concept in the programming language area.
 % that we are facing to formalize and {obtain} the adaptive data analysis program's adaptivity and control its generalization error,
 % with an 
 %
 In Section~\ref{sec:adapt-overview}, we show a running example and give the technical overview of our approach.
 %
 Then the last part in Section~\ref{sec:adapt-outline} shows the outline of \redd{PART \romannum{1}}.
 
 \paragraph{Chapter~\ref{sec:prework}} summarizes the previous works on analyzing the \emph{adaptivity}.
 It contains a short summary of the language designs in Section~\ref{sec:prework-language};
 the adaptivity formalization through program analysis over the trace-based operational semantics in Section~\ref{sec:prework-formalization};
 and the program analysis algorithm that is used to estimate the adaptivity of the data analysis programs 
 in Section~\ref{sec:prework-static}.
%  Their algorithm first introduces the SSA version of the loop language
%  to enable an easier analysis over adaptivity by an easier tracking of dependency relations between variables based on the limitation of direct analysis over the loop language
%  followed with a three-step algorithm, which constructs a data control dependency graph, and add weights to the graph. The adaptivity is estimated by the weight of the path with the highest weight.


 % \paragraph*{Chapter~\ref{sec:adapt-analysis}} presents the new adaptivity analysis framework in following sections.
 
 \paragraph*{Chapter~\ref{sec:adapt-language}} gives our
 {\tt Qwhile} language, which is extended from the standard while language. This chapter includes the syntax in Section~\ref{sec:language-syntax}
 and the trace-based operational semantics in Section~\ref{sec:language-os}.
 % supports more general adadptive data analysis.
 
 \paragraph*{Chapter~\ref{sec:adapt-exe}} shows the semantics-based definition for
 the \emph{adaptivity} of a data analysis program.
 This formal definition is presented through the following sections. 
%

 In Section~\ref{sec:dynamic-datadep}, we define a semantics-based \emph{variable may-dependency} relation to the reason about the dependency relation between every query.
 % through the methodology of semantic data dependency analysis.
%
 In Section~\ref{sec:dynamic-graph}, we first define the quantitative sub-property as the weight for each variable
 based on the \emph{variable may-dependency} relation.
%  Then we encode this quantitative sub-property and with \emph{variable may-dependency} relations  and 
 Then we construct a weighted dependency graph
 namely \emph{semantics-based dependency graph} and encode this quantitative sub-property and with \emph{variable may-dependency} relations on this graph.
 %

 Then in Section~\ref{sec:dynamic-adapt}, we define the \emph{adaptivity} 
%  model in Definition~\ref{def:trace_adapt} as a formal quantitative reachability property.
%  This definition is 
based on the \emph{semantics-based dependency graph} by considering both the \emph{quantitative} and \emph{reachability} sub-properties encoded on it.
% , which captures intuitive \emph{adaptivity} accurately.

 \paragraph*{Chapter~\ref{sec:adapt-static}} introduces a static program adaptivity analysis framework, named {\THESYSTEM} through the following sections.
%  This analysis combines both data dependency analysis with resource cost analysis techniques.
%  It is developed as follows in order to accurately approximate the formal \emph{adaptivity} definition
%  through static program analysis techniques. 

 In Section~\ref{sec:alg_vertexgen},~\ref{sec:alg_weightedgegen} and~\ref{sec:alg_graphgen}, we construct an 
 \emph{estimated dependency} graph for approximating the \emph{semantics-based dependency graph}.
 Specifically, in Section~\ref{sec:alg_edgegen}, we first use an intergrated data-control flow analysis algorithm to soundly
 estimate the \emph{variable may-dependency} relation and build the edges of the estimated graph.
 Then in Section~\ref{sec:alg_weightgen}, we present a reachability-bound analysis algorithm which can soundly
 estimate the quantitative sub-property of the \emph{adaptivity}. Then we use it to construct the weight of the estimated graph.
%  , namely the weight of each variable. 
%  {\THESYSTEM} then uses this algorithm and estimates the \emph{dependency quantity} for \emph{adaptivity}, as well as the weight of this graph.
 At the end, in Section~\ref{sec:alg_graphgen}, we build the 
 \emph{estimate dependency graph} using all these results estimated about.
 In Section~\ref{sec:alg_adaptcompute}, we present an efficient and accurate path search algorithm, which computes the \emph{adaptive} from the estimated graph.
 % , specifically estimating 

 \paragraph{Chapter~\ref{sec:adapt-implementation}} has two sections.
 In Section~\ref{sec:adapt-example}, we give five examples to manually demonstrate our framework.
 Then we show the evaluation results over a benchmark composed of six real-world data analysis algorithms and eleven handcrafted programs based on the code pieces extracted from the C library and six synthesized programs in Section~\ref{sec:adapt-eval}.
 
 \paragraph{Chapter~\ref{sec:adapt-relatedwork}} discusses the related works we have looked in three different areas,
 including the works in semantics-based dependency analysis and formalization area in Section~\ref{sec:relatedwork-adapt-semantics}; the existing static program analysis techniques in Section~\ref{sec:relatedwork-adapt-static}; and the works about generalization error in adaptive data analysis area in Section~\ref{sec:relatedwork-adapt-error}. 

 \item \redd{PART \romannum{2} \quad REACHABILITY BOUND ANALYSIS}

 \paragraph{Chapter~\ref{sec:reachability-intro}} has four parts.

 In Section~\ref{sec:reachability-background}, we introduce the background of the
 \emph{reachability-bound} and developments in analyzing it.
 Then in
 Section~\ref{sec:reachability-motivation}, we discuss
 % and existing 
 the technical limitations and motivations in analyzing the \emph{reachability-bound}.
 Section~\ref{sec:reachability-overview} gives the technical overview of our new approach through
 three running examples.
 % then presents the motivation and existing limitations of this problem and give an overview of the new algorithm in
 % Section~\ref{sec:reachability-motivation} and 
 The structure of this chapter is given 
 in Section~\ref{sec:reachability-outline}.

 \paragraph{Chapter~\ref{sec:psRB-language}} contains the program model including the
 language syntax in Section~\ref{sec:psRB-language-syntax} and the trace-based operational semantics in Section~\ref{sec:psRB-language-semantics}.
 
\paragraph{Chapter~\ref{sec:psRB-definition}} gives the \emph{reachability-bound} definition over our program model.

\paragraph{Chapter~\ref{sec:psRB-analysis-alg}} 
%  presents the {path-sensitive reachability-bound analysis} through the following sections.
%  Section~\ref{sec:psRB-analysis-alg} 
presents the new \emph{path-sensitive reachability-bound} analysis algorithm, named $\kw{psRB}$ through the following subsections.
 
 Section~\ref{sec:psRB-preface} first shows the main steps of this algorithm as a structural section. There are five steps in this algorithm.
 
 In Section~\ref{sec:psRB-progabs}, we give the computation steps for computing the \emph{abstract transition graph}
 for the program.

 Based on this graph, in Section~\ref{sec:psRB-refine}, we refine the multi-paths loop in a program and generates a refined program.

 According to the graph and the refined program, we present the computation steps for computing the ranking function for each edge on the graph and estimates the value of each ranking function  in Section~\ref{sec:psRB-rank}.
 
 In Section~\ref{sec:psRB-psrb} we introduce our \emph{path-sensitive reachability-bound} algorithm, and we give the detailed computations of the \emph{path reachability-bound} in Section~\ref{sec:psRB-pathlocalrb} and Section~\ref{sec:psRB-pathrb},
 and the computation detail of the \emph{loop reachability-bound} in Section~\ref{sec:psRB-looprb}.
%   this algorithm in the next sections.
%  Section~\ref{sec:psRB-pathlocalrb} and Section~\ref{sec:psRB-pathrb} is the algorithm for computing the \emph{path reachability-bound} and  Section~\ref{sec:psRB-looprb} is the  algorithm for computing the \emph{loop reachability-bound}.
%   for each path and loop of this program using the ranking function and the graph.
 

 \paragraph{Chapter~{\ref{sec:psRB-eval}}} has two sections, inlcuding six examples to illustrate the algorithm step-by-step in Section~\ref{sec:psRB-example} firstly,
 and then implementation and evaluation results over 270 programs from four benchmarks sets in Section~\ref{sec:psRB-eval}.
 
  \paragraph{Chapter~\ref{sec:psRB-relatedwork}} is the related works we looked into in program analysis area.

 \item \redd{PART \romannum{3} \quad Epilogue}
 \paragraph{Chapter~\ref{sec:conclusionfuture}}
 is the conclusion of the works in \redd{PART \romannum{1}} and \redd{PART \romannum{2}}, and the future plans based on the limitations through the following sections.
%  Section~\ref{sec:conclusionfuture-adapt} and Section~\ref{sec:conclusionfuture-psRB}.

In Section~\ref{sec:conclusion-adapt}, 
%  summarizes the program analysis framework for adaptive data analysis in \redd{PART \romannum{1}}, and 
we conclude the contribution and weakness of our program analysis framework for adaptive data analysis. Then in Section~\ref{sec:future-adapt}, we propose future work on extending and improving our existing work.

 Section~\ref{sec:conclusion-psRB} summarizes the contributions and limitations of our path-sensitive reachability-bound computation algorithm presented in \redd{PART \romannum{2}}. Then Section~\ref{sec:future-psRB} proposes the future works, specifically the improvement of existing works and the possible applications on other fields.
 
%  Section~\ref{sec:future-cost} discusses a possible future direction for studying general quantitative reachability properties, i.e., {the program's non-monotonic quantitative property}.

 \item \redd{PART \romannum{4} \quad APPENDIX}
 The Appendix contains the complete program code for examples, and proofs for lemmas and theorems in \redd{PART \romannum{1}} and \redd{PART \romannum{2}}.
\end{enumerate}

% \textbf{Previous Published Materials}